\documentclass[11pt,a4paper]{article}
\usepackage[margin=1in]{geometry}
\usepackage[T1]{fontenc}
\usepackage[utf8]{inputenc}
\usepackage{lmodern,microtype,amsmath,amssymb,amsthm,mathtools}
\usepackage[colorlinks=true,linkcolor=blue,urlcolor=blue]{hyperref}
\hypersetup{pdftitle={Learning parity from selected examples under a globally separated neural kernel},pdfauthor={Dmitry Vaintrob; Kaarel Hanni; OpenAI Codex}}
\setlength{\emergencystretch}{2em}
\newtheorem{theorem}{Theorem}
\newtheorem{lemma}[theorem]{Lemma}
\newtheorem{corollary}[theorem]{Corollary}
\theoremstyle{remark}\newtheorem{remark}[theorem]{Remark}
\newcommand{\E}{\mathbb E}\newcommand{\Pp}{\mathbb P}
\newcommand{\normK}[1]{\lVert #1\rVert_K}
\title{Learning parity from selected examples\\under a globally separated neural kernel}
\author{Dmitry Vaintrob \and Kaarel H\"anni \and OpenAI Codex}
\date{October 2026}
\begin{document}
\maketitle
\begin{abstract}
We construct selected training sets on which a globally margin-conditioned
neural Gaussian prior learns sparse parity using $d^{O(\sqrt{k})}$ examples.
For any fixed error tolerance, every minimum-energy compatible field has
that error or less, and all such fields have the same signs. The construction
uses alternating interpolation on central Hamming layers and a general
sparsification lemma for convex margin minimizers. It also produces actual
global minimizers with very large amplitudes at correctly classified inputs.
These are deliberately selected examples, not independent uniform samples.
The result does not settle the proposed $\exp\{\Omega(k\log(ed/k))\}$ lower
bound for uniform random training; it shows that typical sampling is essential
to such a bound. It suffices that $d\ge(C_Lk)^{C_Lk}$.
\end{abstract}

\section{Model and result}
Let $\mathcal X_d=\{-1,1\}^d$ have uniform probability measure. Fix
$L\ge2$, define $\Psi(t)=2\arcsin(t)/\pi$, and put
\[
 F=\Psi^{\circ(L-1)},\qquad
 K_d(x,y)=F\!\left(\frac{1+x\cdot y}{d+1}\right).
\]
This covariance is strictly positive definite and has diagonal one. For a
real field $z$ on the cube, its energy is
\[
 \normK z^2=z^{\mathsf T}K_d^{-1}z.
\]
The factor $1/2$ in Gaussian rate functions does not affect the minimizers.
The unit margin entails no loss of generality: multiplying fields by a
fixed margin $a>0$ multiplies all energies by $a^2$ and preserves signs.

The target is $f(x)=\prod_{i=1}^k x_i$. Given a finite set of training
inputs $T$, consider
\begin{equation}\label{eq:global}
 \min_z\normK z^2\quad\text{subject to}\quad
 |z(x)|\ge1\ (x\in\mathcal X_d),\qquad f(x)z(x)\ge1\ (x\in T).
\end{equation}
This is a global separation constraint at \emph{every} input. A minimizer's
error is $\Pp_x\{f(x)z(x)<0\}$.

\begin{theorem}[Selected examples]\label{thm:main}
Fix $L\ge2$ and $\delta\in(0,1/2)$. For every sufficiently large odd $k$
and every $d\ge(C_Lk)^{C_Lk}$, there is a parity-labeled training set
$T\subseteq\mathcal X_d$ with
\[
 |T|\le d^{\,2\sqrt{k/\delta}+9}
\]
for which all minimizers of~\eqref{eq:global} have the same sign function
and error at most $\delta$. At a correctly classified input, every such
minimizer has amplitude at least
\[
 \exp\{c_\delta\sqrt{k}\log k\}.
\]
The positive constant $c_\delta$ is independent of $k,d$; $C_L$ depends
only on $L$, and the threshold on $k$ may depend on $\delta$.
\end{theorem}

The explicit dimension condition permits growing $k$ up to a constant
multiple of $\log d/\log\log d$. The training set is chosen
using a contact measure of a convex margin problem. It is not distributed
as independent uniform examples. Nevertheless, for fixed $k$ the displayed
number of examples is a vanishing fraction of the $2^d$ inputs. For growing
$k$, its exponent in $d$ is sublinear in $k$.

\begin{corollary}[Actual label probabilities]\label{cor:gp}
For the selected set in Theorem~\ref{thm:main}, let $Z_t$ be Gaussian with
covariance $tK_d$, conditioned on $|Z_t(x)|>1$ at every input and on the
correct training signs. As $t\downarrow0$, the posterior probability of
the common minimizing sign tends to one at every input. Consequently,
\[
 \lim_{t\downarrow0}\Pp_x\left\{
   \Pp\{f(x)Z_t(x)>0\mid\text{conditioning}\}\le2/3\right\}\le\delta.
\]
\end{corollary}

The statement concerns sign probabilities, not the Gaussian posterior
mean. The variational problem~\eqref{eq:global} is also the one that arises
in the width-first subcritical neural regime $0<\gamma<1$. The proofs below
are entirely finite dimensional and use its explicit kernel; they do not
establish or require a neural width-limit theorem.

\section{Fixed-degree eigenvalues}
For $S\subseteq[d]$, write $\chi_S(x)=\prod_{i\in S}x_i$. Let
$\lambda_j(d)$ denote the eigenvalue on degree-$j$ characters of the
\emph{probability-normalized} kernel operator $2^{-d}K_d$. Thus
\begin{equation}\label{eq:walshenergy}
 \normK z^2=\sum_{S\subseteq[d]}
       \frac{\widehat z(S)^2}{\lambda_{|S|}(d)}.
\end{equation}

\begin{lemma}\label{lem:spectrum}
The matrix $K_d$ is strictly positive definite. For every fixed positive
odd integer $j$,
\begin{equation}\label{eq:spectralasymptotic}
 \lambda_j(d)=F^{(j)}(0)(d+1)^{-j}+O_j(d^{-j-1}),
 \qquad F^{(j)}(0)>0.
\end{equation}
The implicit constant may depend on $L,j$. Moreover, for a constant $C_L$
and $B_k=(C_Lk)^{C_Lk}$, for every $d\ge k$ and odd $1\le j\le k$,
\begin{equation}\label{eq:uniformspectrum}
 B_k^{-1}d^{-j}\le\lambda_j(d)\le B_kd^{-j}.
\end{equation}
\end{lemma}
\begin{proof}
The odd Taylor series of $\Psi$ has strictly positive coefficients at
every odd degree. Composition preserves this property. Expand
$F(t)=\sum_{m\text{ odd}}c_mt^m$. The polynomial kernel
$((1+x\cdot y)/(d+1))^m$ has nonnegative Walsh eigenvalues. For any
character of reduced degree $j$, the term $m=j$ if $j$ is odd, and the
term $m=j+1$ if $j$ is even, gives a strictly positive coefficient:
in the latter case use the additional constant coordinate once. This
proves strict positive definiteness.

For odd $j$, write
\[
 \lambda_j(d)=\E\left[
 F\!\left(\frac{1+\sum_{i=1}^d X_i}{d+1}\right)
 \prod_{i=1}^jX_i\right]
\]
with independent uniform signs $X_i$. Taylor-expand $F$ through degree
$2j+1$ on $[-1/2,1/2]$. The complementary event has exponentially small
probability; both $F$ and the fixed Taylor polynomial are bounded on
$[-1,1]$. Oddness makes the remainder $O_j(|t|^{2j+3})$, so it contributes
$O_j(d^{-j-3/2})$, by the fixed moments of a Rademacher sum. Terms of
degree less than $j$ vanish. The degree-$j$ term contributes
$c_jj!/(d+1)^j$. For a fixed $m>j$, the signed moment of
$(1+\sum X_i)^m$ is $O_j(d^{(m-j)/2})$: each of the specified $j$
coordinates must occur an odd number of times, while other coordinates
must occur an even number of times. Since the next odd $m$ is $j+2$,
all remaining Taylor terms are $O_j(d^{-j-1})$.

For the uniform bounds, every Taylor coefficient of $F$ is at most one,
because they are nonnegative and sum to $F(1)=1$. At odd degrees,
$c_j(F)\ge(2/\pi)^{L-2}c_j(\Psi)\ge c_Lj^{-3/2}$, by keeping the
linear term in each outer composition. The degree-$j$ term alone proves
the lower bound in~\eqref{eq:uniformspectrum}, absorbing $2^j$ and the
fixed constants into $B_k$.

For the upper bound, Taylor-expand through degree $2k+1$. For each odd
$m\le2k+1$ with $m\ge j$, the preceding signed-moment count is bounded
by $(Ck)^{Ck}d^{(m-j)/2}$: at most $(m-j)/2$ distinct unspecified
coordinates can occur, and all multiplicities and orderings cost at most
$(Ck)^{Ck}$. Summing the finitely many terms gives $B_kd^{-j}$. On
$|t|\le1/2$, oddness and the coefficient bound give remainder at most
$2|t|^{2k+3}$; its expectation is at most
$(Ck)^{k+3/2}d^{-k-3/2}$. On the complementary event the error is at
most $Ck$, with probability at most $Ce^{-cd}$. Finally,
$\sup_{d>0}d^ke^{-cd}\le(Ck)^k$. These errors are absorbed into
$B_kd^{-j}$, after increasing $C_L$.
\end{proof}

\section{A globally separated convex minimizer}
Choose $r$ to be the smallest integer satisfying
\begin{equation}\label{eq:rchoice}
 r\ge\sqrt{k/\delta},\qquad r\equiv k\pmod4.
\end{equation}
For sufficiently large odd $k$, $r$ is odd and $r<k$. Set
$h_0=(k-r)/2$, which is even. If $H(x)$ is the number of negative active
coordinates, let
\[
 C=\{x:h_0\le H(x)\le h_0+r\},
\]
including every nuisance-coordinate assignment. Chebyshev's inequality,
using $\operatorname{Var}(H)=k/4$, gives
\begin{equation}\label{eq:mass}
 \Pp(C^c)=\Pp\{|H-k/2|>r/2\}\le k/r^2\le\delta.
\end{equation}
Initially impose only the convex constraints on $C$:
\begin{equation}\label{eq:convex}
 z_d=\arg\min_z\{\normK z^2:f(x)z(x)\ge1\text{ for }x\in C\}.
\end{equation}

Let $p_r$ be the degree-$r$ polynomial interpolating
$p_r(i)=(-1)^i$, $0\le i\le r$, and define
\begin{equation}\label{eq:interpolant}
 P(x)=(-1)^{h_0}p_r(H(x)-h_0).
\end{equation}
It has degree $r$ on the active cube and training margins exactly one.

\begin{lemma}\label{lem:interpolant}
There is a constant $C_L$ such that, for $B_k=(C_Lk)^{C_Lk}$ and
all $d\ge k$,
\[
 \max_x|z_d(x)-P(x)|\le B_k/d,\qquad
 E_d:=\normK{z_d}^2\le B_kd^r.
\]
Consequently, for $d\ge2B_k$, $z_d$ is globally separated, is the unique
minimizer of~\eqref{eq:global} with $T=C$, and has error at most $\delta$.
\end{lemma}
\begin{proof}
Strict convexity gives uniqueness in~\eqref{eq:convex}. Invariance under
nuisance flips makes $z_d$ independent of nuisance coordinates; active
permutations make it a function of $H$. Since $k$ is odd and the layer
interval is symmetric, the operation $z(x)\mapsto-z(-x_{[k]},x_{[d]\setminus[k]})$
also preserves the problem. Thus $z_d$ is odd under $H\mapsto k-H$ and
has only odd active Walsh levels. The same symmetry holds for $P$, by
uniqueness of interpolation.

In this proof, $B$ denotes a quantity bounded by $(C_Lk)^{C_Lk}$; the
constant $C_L$ can increase a fixed number of times. Products, fixed powers,
and factors such as $k!$ or $2^k$ are absorbed into this convention.
Use the fixed orthonormal symmetric active Walsh basis
\[
 \phi_j=\binom kj^{-1/2}\sum_{|S|=j,\ S\subseteq[k]}\chi_S,
 \qquad j\text{ odd}.
\]
These are polynomials of degree $j$ in $H$. Their leading coefficients
and $\max_H|P(H)|$ are bounded by $B$: for the latter, use
$p_r(t)=\sum_{j=0}^r(-2)^j\binom tj$ and $|t|\le k$.
Write $c_j(z)$ for the coefficient of $\phi_j$.
By~\eqref{eq:uniformspectrum}, the comparator $P$ gives
\[
 \normK{z_d}^2\le\normK P^2\le Bd^r.
\]
Let $u_d$ be the degree-at-most-$r$ part of $z_d$, and $v_d=z_d-u_d$.
Since the first possible higher degree is $r+2$, the same spectral bounds
and $\|v_d\|_\infty\le2^{k/2}\|v_d\|_2$ give
\begin{equation}\label{eq:highpart}
 \|v_d\|_\infty\le B/d.
\end{equation}

A direct coefficient comparison is stronger than merely bounding the
energy of $u_d$. Namely,
\begin{align*}
 |c_r(z_d)|^2
 &\le\lambda_r(d)\normK P^2\\
 &=|c_r(P)|^2+
   \sum_{j<r,\ j\text{ odd}}
        \frac{\lambda_r(d)}{\lambda_j(d)}|c_j(P)|^2
 \le |c_r(P)|^2+B/d^2.
\end{align*}
Consequently
\begin{equation}\label{eq:leadbound}
 |\operatorname{lead}(u_d)|\le|\operatorname{lead}(P)|+B/d.
\end{equation}
Put $y_{i,d}=(-1)^{h_0+i}z_d(h_0+i)\ge1$. Interpolation at consecutive
nodes and~\eqref{eq:highpart} give
\[
 \operatorname{lead}(u_d)=
 (-1)^{h_0+r}\sum_{i=0}^r\frac{y_{i,d}}{i!(r-i)!}+e_d,
 \qquad |e_d|\le B/d.
\]
All these weights are positive, and
$|\operatorname{lead}(P)|=\sum_i1/[i!(r-i)!]$. Thus
\eqref{eq:leadbound} implies
\[
 \sum_{i=0}^r\frac{y_{i,d}-1}{i!(r-i)!}\le B/d.
\]
Each summand is nonnegative and each weight is at least $1/r!$, so
$0\le y_{i,d}-1\le B/d$ for every $i$. The Lagrange interpolation
operator from these $r+1$ consecutive nodes to all $H\in\{0,\ldots,k\}$
has norm at most $(r+1)k^r\le B$: each denominator is a nonzero integer
and each numerator has $r$ factors of magnitude at most $k$.
Combining this with~\eqref{eq:highpart} proves
$\|z_d-P\|_\infty\le B/d$.

There is a root of $p_r$ in each interval $(i,i+1)$, $0\le i<r$;
these exhaust its roots. Its absolute value increases as one moves
outward from $[0,r]$. Newton interpolation gives
$p_r(r+1)=1-2^{r+1}$ for odd $r$, and reflection gives the same
magnitude at $t=-1$. Consequently $|P|\ge2^{r+1}-1$ on every integer
layer outside $C$. If $d\ge2B_k$, the uniform difference is at most
$1/2$, so $|z_d|>1$ outside $C$; on $C$, its prescribed margins are
already at least one. Thus $z_d$ is feasible for the global problem.
Since~\eqref{eq:convex} is a relaxation, it is the unique global minimizer
as well. Equation~\eqref{eq:mass} bounds its error.
\end{proof}

\begin{lemma}[A large correct amplitude]\label{lem:amplitude}
At the active input with $H=0$, the target and $P$ are positive and
\[
 P(0)\ge\left(\frac{k-r}{r}\right)^r.
\]
Under the dimension condition $d\ge2B_k$ from
Lemma~\ref{lem:interpolant}, $z_d(0)\ge P(0)/2$.
\end{lemma}
\begin{proof}
The leading coefficient of $p_r$ is $(-2)^r/r!$. All roots of $P$,
viewed as a polynomial in $H$, lie between $h_0$ and $h_0+r$. Since
$h_0$ is even and $r$ is odd, $P(0)>0$. Therefore
\[
 P(0)\ge\frac{2^r h_0^r}{r!}
       \ge\left(\frac{2h_0}{r}\right)^r.
\]
The final claim follows from Lemma~\ref{lem:interpolant}.
\end{proof}

\section{Sparsifying the selected examples}
The next lemma applies to any finite positive definite covariance. It
controls all global minimizers, even when the globally separated feasible
set is nonconvex.

\begin{lemma}[Sparse teaching by contact sampling]\label{lem:teaching}
Let $K(x,x)\le1$. Suppose $z_*$ is the convex minimum-energy solution of
unit-margin constraints on a set $C$ and also satisfies $|z_*|\ge1$
everywhere. Write $E=\normK{z_*}^2$. For every positive integer $n$,
there is a multiset $T$ of $n$ contact points in $C$ such that every
minimum-energy globally separated field $z_T$ consistent with their
prescribed labels satisfies
\begin{equation}\label{eq:teachingdistance}
 \normK{z_T-z_*}\le\frac{2E}{\sqrt n},\qquad
 \max_x|z_T(x)-z_*(x)|\le\frac{2E}{\sqrt n}.
\end{equation}
Repeated points may be discarded.
\end{lemma}
\begin{proof}
Convex duality and complementary slackness give multipliers $\mu_x\ge0$
supported on contacts, with
\[
 z_*=Kh,\qquad h_x=f(x)\mu_x,\qquad
 \sum_x\mu_x=E,\qquad h^{\mathsf T}Kh=E.
\]
Draw $X_1,\ldots,X_n$ independently with law $\mu/E$, and set
\[
 h_T=\frac En\sum_{i=1}^nf(X_i)e_{X_i}.
\]
It is a feasible dual vector for the convex margin problem on $T$, with
objective $D_T=2E-h_T^{\mathsf T}Kh_T$. Independence and the diagonal
bound imply
\[
 \E[h_T^{\mathsf T}Kh_T]\le(1-1/n)E+E^2/n.
\]
Some realization therefore has $D_T\ge E-E^2/n$. Let $v_T$ be the
convex minimum-norm field satisfying just the training margins on this
realization, and let $E_T=\normK{v_T}^2$. Weak duality gives
$E_T\ge E-E^2/n$.

Both $z_*$ and every globally separated minimizer $z_T$ are feasible for
this convex problem. Also $\normK{z_T}^2\le E$, because $z_*$ is a
feasible globally separated comparator. The Hilbert projection inequality
at $v_T$ consequently gives
\[
 \normK{z_*-v_T}^2\le E-E_T,\qquad
 \normK{z_T-v_T}^2\le E-E_T.
\]
The triangle inequality proves the first claim in~\eqref{eq:teachingdistance}.
The reproducing property and $K(x,x)\le1$ prove the second.
\end{proof}

\begin{proof}[Proof of Theorem~\ref{thm:main}]
Increase $C_L$ in the theorem so that
Lemma~\ref{lem:interpolant} applies for every permitted $d$. Use Lemma~\ref{lem:teaching} with
$z_*=z_d$ and $n=\lceil16E_d^2\rceil$. Every global minimizer on the
resulting selected set is within $1/2$ of $z_d$ at every input. Since
$|z_d|\ge1$, all have its sign function, hence error at most $\delta$.
At $H=0$ their amplitude is at least $P(0)/2-1/2$, which is
$\exp\{c_\delta\sqrt{k}\log k\}$ for sufficiently large $k$ by
Lemma~\ref{lem:amplitude} and~\eqref{eq:rchoice}.

The energy bound gives $n\le 17B_k^2d^{2r}$. Increasing $C_L$ once
more ensures that the dimension condition absorbs $17B_k^2$ into one
power of $d$. Thus $|T|\le n\le d^{2r+1}$.
Since $r<\sqrt{k/\delta}+4$, this is bounded by the power stated in
the theorem.
\end{proof}

\begin{proof}[Proof of Corollary~\ref{cor:gp}]
There are finitely many sign orthants. The closed feasible set in each
compatible orthant has a minimum Gaussian rate $\normK z^2/2$, by
coercivity. All global minimizers have the common sign from
Theorem~\ref{thm:main}, so every other compatible orthant has a strictly
larger minimum rate. Taking the minimum of finitely many positive gaps
still gives a positive gap. Strict inequalities in the conditioning do
not change the rate infimum: scaling a closed feasible field by
$1+\varepsilon$ makes every absolute margin strict. Elementary
finite-dimensional Gaussian Laplace bounds therefore make the total
posterior probability of the other orthants tend to zero. The cube is
finite, so this gives the stated pointwise probability and confidence
conclusion.
\end{proof}

\section{Why this does not answer the random-sample problem}
The contact sampling in Lemma~\ref{lem:teaching} selects only inputs where
the dense convex minimizer already agrees with parity. Uniform sampling
also produces examples in its error region. Then that field is no longer
feasible, and the upper energy bound $\normK{z_T}^2\le E$ used in the
sparsification proof fails. Reweighting uniform observations can approximate
the contact dual, but it cannot discard those conflicting constraints.

Accordingly, the theorem rules out a $d^{\Omega(k)}$ requirement that is
claimed for \emph{every} parity-consistent training set. It neither proves
nor disproves the natural lower bound for independent uniform examples.
It also rules out a correct-amplitude ceiling derived solely from global
optimality and the contact equations. Any such control in the uniform
sampling problem must use its probabilistic structure.
\end{document}
